problem:  
Let \(X\) be a smooth complex Fano manifold and \( \omega_0 \in 2\pi c_1(X) \).  
For each initial Kähler potential \( \varphi_0 \) with \( \omega_{\varphi_0} = \omega_0 + i\partial\bar{\partial}\varphi_0 \), consider the **normalized Kähler–Ricci flow**
\[
\frac{\partial \varphi_t}{\partial t} = \log \frac{ \omega_{\varphi_t}^n }{ e^{h_{\omega_0}-\varphi_t}\omega_0^n }, 
\quad \varphi_t|_{t=0} = \varphi_0.
\]
If \(X\) is \(K\)-stable, the flow is expected to converge to a Kähler–Einstein metric.  
When \(X\) is \(K\)-unstable, the Ding functional \( \mathcal{D}(\varphi_t) \) decreases indefinitely along the flow.

**Question:**  
Assume \(X\) is \(K\)-unstable.  
Does the Kähler–Ricci flow, after natural time and potential normalization, asymptotically align with a *canonical destabilizing geodesic ray* in the \(d_1\)-metric completion of the space of Kähler potentials?  
More precisely:  

- Does there exist a geodesic ray \( \Phi(t) \in \mathcal{E}^1(X,\omega_0) \) such that the \(d_1\)-distance between the flow trajectory \( \varphi_t \) and the truncated segment \( \Phi(t) \) remains sublinear in \(t\):  
  \[
  \lim_{t\to\infty} \frac{d_1(\varphi_t, \Phi(t))}{t} = 0,
  \]
  and \( \Phi \) achieves the minimal Ding slope among all finite-energy rays?  
- If such a geodesic ray exists, is it unique up to translation and scaling, and is its associated non‑Archimedean metric the *optimal destabilizer* in the sense of Blum–Jonsson?  

Why is it a "good" problem:  
This problem connects **evolutionary flow dynamics** (the long‑time behavior of the Kähler–Ricci flow) with **variational instability theory** (Ding functional slopes and non‑Archimedean metrics).  
It asks whether a purely analytic flow—in the absence of stability—**selects a canonical destabilizing direction** among all geodesic rays.  

This is a sharply formulated and currently unanswered question: even in unstable settings, the asymptotic geometry of the Kähler–Ricci flow is poorly understood.  
A positive result would unify stability theory and parabolic analysis by showing that the Ricci flow “detects” the non‑Archimedean optimal degeneration dynamically.  
A counterexample would demonstrate a genuine gap between the dynamical and variational pictures of instability.  

The problem bridges deep areas—geometric analysis (Ricci flow asymptotics), pluripotential geometry (\(d_1\)-metric and geodesic rays), and algebro‑geometric stability theory—through a clear, concrete formulation.  
It is narrow, precise, and potentially transformative, hence a genuinely “good” research-level problem.